\documentclass[11pt, a4paper]{article}

\usepackage[margin=2.5cm]{geometry}
\usepackage{amsmath, amssymb}
\usepackage{booktabs}
\usepackage{hyperref}
\usepackage{xcolor}
\usepackage{enumitem}
\usepackage{titlesec}
\usepackage{fancyhdr}

\hypersetup{colorlinks=true, linkcolor=blue!70!black, urlcolor=blue!70!black}

\pagestyle{fancy}
\fancyhf{}
\rhead{TFE Structural Physics --- Confidential}
\lhead{Structural Exit Assessment Specification}
\rfoot{Page \thepage}

\titleformat{\section}{\large\bfseries}{}{0em}{}[\titlerule]
\titleformat{\subsection}{\normalsize\bfseries}{}{0em}{}

\begin{document}

\begin{titlepage}
  \centering
  \vspace*{2cm}
  {\Huge\bfseries Structural Exit Assessment\\[0.3em]
  Specification\par}
  \vspace{0.5cm}
  {\large\itshape Full-Tuple Deterioration Analysis\\
  for Position Hold/Cut Decisions\par}
  \vspace{2cm}
  {\large Tao Financial Engine --- Internal Specification\par}
  \vspace{0.5cm}
  {\normalsize Version 1.0 \quad|\quad May 2026\par}
  \vspace{0.5cm}
  {\normalsize Status: \textbf{Approved Design --- Implementation Pending}\par}
  \vspace{3cm}
  \begin{abstract}
    \noindent
    This specification defines the structural exit assessment for held positions
    in the Tao Financial Engine. When a position enters the intermediate loss zone
    (between 0\% and $-10$\%), the system assesses whether the structural energy
    that justified the entry has dissipated or is merely paused. The assessment
    reads the full coupled kernel tuple as a single geometric object, using three
    inputs: the horse (the stock's own structural history with this shape), the
    herd (structural peers assessed individually then read collectively), and
    match confidence (distance to closest precedent). The $-10$\% catastrophic
    floor remains unconditionally active regardless of assessment outcome.
  \end{abstract}
\end{titlepage}

\tableofcontents
\newpage

% ─────────────────────────────────────────────────────────────────────────────
\section{Motivation}

\subsection{The Exit Gap}

The current exit logic checks individual structural fields:
\begin{itemize}
  \item EXIT-A: $S_{UF} \geq 0.75$ (acceleration complete)
  \item EXIT-B: $D_k \neq 1$ (directional collapse)
  \item EXIT-F: price loss $\geq 10\%$ (catastrophic floor)
  \item EXIT-R9: no losing exit before 7 calendar days
\end{itemize}

Between 0\% and $-10$\%, positions sit with no intermediate structural check.
A stock can be down 7\% with its structural thesis completely broken across
every kernel field, and the system will hold it for the full 7-day window
(or until $-10$\%) because no single field has crossed its individual threshold.

\subsection{The Principle}

The kernel produces a coupled tuple:
\[
  \boldsymbol{\psi} = (D_k, M_k, B_k, R_{rev,k}, U^*_k, C_k, P_k, R_k)
\]

This tuple describes a single geometric structural state. It is not eight
independent signals. The coupling carries information that individual field
checks destroy.

A position where $D_k = 0$ but $M_k > 0$, $B_k$ holding, $S_{UF}$ stable
is a fundamentally different structural state than one where $D_k = 0$ and
all other fields are also deteriorating --- even though a single-field exit
check on $D_k$ would treat them identically.

\textbf{The exit assessment must read the full tuple as one geometric object.}
This is the same principle as entry selection (KERNEL\_PHILOSOPHY.md) applied
to exit decisions.

% ─────────────────────────────────────────────────────────────────────────────
\section{The Diamond Analogy}

A gemologist does not evaluate a diamond by checking clarity separately,
then carat, then color, then voting. The gemologist reads the stone as a
whole: ``VS2 H round brilliant'' is a single classification that captures
the full profile. The individual characteristics inform the grade, but the
grade is the thing.

The structural exit assessment works the same way. The full tuple vector
$\boldsymbol{\psi}$ is the diamond. We find historical diamonds with the
same cut (nearest neighbors in the full tuple space) and ask: what happened
to those diamonds?

We do \textbf{not} decompose $\boldsymbol{\psi}$ into individual fields
and narrate each one. The distance metric reads all facets simultaneously.
The output is: ``this shape historically leads to recovery 62\% of the time
for this stock, and 70\% of the time for its structural peers.'' Not:
``$D_k$ says go but $M_k$ says slow down.''

% ─────────────────────────────────────────────────────────────────────────────
\section{Assessment Method}

\subsection{When It Runs}

The assessment triggers when a held position enters the intermediate loss
zone: unrealized loss exceeds a threshold (initially $-2\%$) but has not
yet reached the $-10\%$ catastrophic floor. It does \textbf{not} run on
every sentinel cycle --- only when deterioration is detected.

\subsection{Three Inputs}

\subsubsection{Input 1: The Horse (Self-History)}

For the stock in question, retrieve its full structural history from
\texttt{runtime\_decisions\_history} (production) or the governed states
parquet (backtest).

Compute the Euclidean distance from the current tuple $\boldsymbol{\psi}_{\text{now}}$
to every historical tuple $\boldsymbol{\psi}_i$ for this stock:
\[
  d_i = \left\| \frac{\boldsymbol{\psi}_{\text{now}} - \boldsymbol{\psi}_i}
  {\boldsymbol{r}} \right\|_2
\]
where $\boldsymbol{r}$ is the per-field range vector computed from this
stock's own history (normalizing so each field contributes equally to the
distance, preventing high-magnitude fields from dominating).

Take the $N$ closest matches (initially $N = 30$). For each match, look up
what the price did at 10, 20, and 30 days forward. Report:
\begin{itemize}
  \item 20-day win rate across the $N$ matches
  \item Average 20-day return
  \item Distance of the closest match (match confidence)
\end{itemize}

\subsubsection{Input 2: The Herd (Structural Peers)}

Find the stock's \emph{structural peers}: stocks whose structural personality
profile (mean tuple over their full history) is most similar.

For each of the top 20 structural peers:
\begin{enumerate}
  \item Run that peer's \textbf{own} full history through the same distance
        computation against the target tuple
  \item Take its top 20 closest self-matches
  \item Compute its 20-day WR from those matches
  \item Classify: BULL ($\geq 60\%$), MIXED (40--60\%), BEAR ($\leq 40\%$)
\end{enumerate}

Read the collective verdict: count of BULL / BEAR / MIXED across all 20 peers.

This is \textbf{not} a single nearest-neighbor lookup across the universe.
Each peer is assessed individually against its own history, then the verdicts
are read collectively. This captures swarm behavior: are structurally similar
stocks recovering from this shape, or deteriorating?

\subsubsection{Input 3: Match Confidence}

The Euclidean distance of the closest self-match. Tight distance means the
stock has been in this exact geometric shape before and the precedent is
strong. Wide distance means unfamiliar territory and the precedent is weak.

\subsection{Epoch Weighting}

The assessment should be weighted by the current epoch state. If the macro
environment is hostile for this stock's sector (epoch pressure $< -0.5$),
discount the bull case. If the macro environment is favorable, the bull
case is reinforced.

This is not a veto --- it is a weighting factor on the structural precedent.

% ─────────────────────────────────────────────────────────────────────────────
\section{Decision Rules}

\subsection{Structural Energy Dissipated --- CUT}

All three of the following:
\begin{enumerate}
  \item The horse's full-tuple self-match 20-day WR $\leq 50\%$
  \item The herd's collective verdict: BULL count $\leq$ BEAR $+$ MIXED
  \item The closest self-match distance is wide (above the stock's own
        historical median match distance)
\end{enumerate}

\textbf{Action:} Exit the position. Do not wait for $-10$\% floor.
Do not wait for 7-day minimum hold to expire. The structural thesis
is dead.

\subsection{Structural Energy Paused --- HOLD}

All three of the following:
\begin{enumerate}
  \item The horse's full-tuple self-match 20-day WR $> 55\%$
  \item The herd's collective verdict: BULL count $> 2 \times$ BEAR count
  \item The closest self-match distance is tight (below the stock's own
        historical median match distance)
\end{enumerate}

\textbf{Action:} Hold through the intermediate loss. The structural energy
is stored but delayed (false start / magma pattern). Extend the observation
window. Watch for the resumption signal.

\textbf{The $-10$\% catastrophic floor remains unconditionally active.}
The hold decision says ``don't cut early on intermediate noise.'' The floor
says ``if the structural read was wrong, protect capital.'' These are not
contradictory --- they operate at different scales. Even when all instruments
say eruption is likely, you still have an evacuation line.

\subsection{Ambiguous --- DEFAULT}

If neither CUT nor HOLD criteria are met:
\begin{itemize}
  \item Fall through to existing exit logic (EXIT-B, EXIT-C, EXIT-R9)
  \item No additional action
\end{itemize}

% ─────────────────────────────────────────────────────────────────────────────
\section{Data Requirements}

\begin{table}[h!]
\centering
\begin{tabular}{lll}
\toprule
\textbf{Data} & \textbf{Source} & \textbf{Notes} \\
\midrule
Stock's own tuple history & \texttt{runtime\_decisions\_history} & Per-bar, all fields \\
Forward returns & Polygon API / cached bars & 10d, 20d, 30d \\
Structural personality profiles & Precomputed weekly batch & Mean/std of tuple per stock \\
Epoch state & \texttt{g32\_state.json} & Current sector pressures \\
Current tuple & \texttt{runtime\_decisions\_latest} & \texttt{snapshot\_row\_json} \\
\bottomrule
\end{tabular}
\end{table}

% ─────────────────────────────────────────────────────────────────────────────
\section{Relationship to Other Specifications}

\subsection{3WA (Three-Wave Alignment)}

The 3WA model (\texttt{docs/structural\_wave\_alignment\_spec.tex}) addresses
\emph{entry} quality --- identifying high-confidence crystallisation events.
This specification addresses \emph{exit} quality --- assessing whether a
held position's structural thesis remains intact.

They are complementary. A 3WA-tagged entry that enters the intermediate loss
zone would be assessed by this exit framework like any other position.

\subsection{Per-Stock Relative Thresholds (Step 2)}

The structural personality profiles computed for per-stock thresholds (species
classification, $\bar{\delta}_i$, $\sigma_i$) are the same profiles used to
identify structural peers in the herd assessment. The data infrastructure
is shared.

\subsection{EXIT-R9 (7-Day Minimum Hold)}

EXIT-R9 blocks losing exits for 7 calendar days. This specification can
override that hold when the CUT criteria are met --- the structural thesis
is dead and waiting 7 days only deepens the loss.

When the HOLD criteria are met, EXIT-R9's hold is reinforced --- the
structural precedent supports waiting.

% ─────────────────────────────────────────────────────────────────────────────
\section{Implementation Notes}

\begin{enumerate}
  \item \textbf{Performance:} The horse lookup queries one stock's history
        ($\sim$1,200 rows for 5-year established stocks). The herd lookup
        queries 20 peers $\times$ $\sim$1,200 rows each. Total: $\sim$25,000
        rows. This runs only on deteriorating positions, not every cycle.
        Expected load: 1--5 positions per assessment, a few times per week.
  \item \textbf{Caching:} Structural personality profiles should be cached
        and refreshed weekly. Peer lists are stable over days/weeks.
  \item \textbf{Forward returns:} In production, forward returns are not
        available (we don't know the future). The horse/herd WR is computed
        from \emph{historical} forward returns at historical match points ---
        ``when this stock was in this shape 6 months ago, what happened
        over the next 20 days?'' This is backward-looking analysis, not
        forward prediction.
  \item \textbf{Production vs quarantine data:} Production kernel uses
        log-normalized L0. Quarantine uses raw. The tuple shapes should be
        similar but absolute values differ. The assessment should use
        production's own \texttt{runtime\_decisions\_history}, not quarantine
        data.
\end{enumerate}

% ─────────────────────────────────────────────────────────────────────────────
\section{Validation}

Before deployment:
\begin{enumerate}
  \item Run the assessment on all closed positions in the trade ledger
        that hit $-5\%$ or worse. Would CUT have fired earlier? Would HOLD
        have correctly identified false starts?
  \item Compare to actual outcomes: of positions that recovered, how many
        would HOLD have correctly identified? Of positions that continued
        to $-10\%$, how many would CUT have caught earlier?
  \item Shadow-run alongside existing exits for 2 weeks before promoting
        to active.
\end{enumerate}

\vspace{2em}
\noindent\textit{This specification was developed from a structural assessment
session on May 27, 2026, using BSM and VIRT as live case studies. The
horse/herd/topology framework and the diamond analogy are from Joseph.
The full-tuple-as-single-object constraint is a core architectural
principle documented in KERNEL\_PHILOSOPHY.md.}

\end{document}
